% ============================================================================
% Fichier  : partie6/P06_C15_recherche_auditologie.tex
% Titre    : Ouvrir sur la Recherche : l'Auditologie
% Auteur   : MINKA MI NGUIDJOÏ Thierry Emmanuel
% Version  : 0.1 -- novembre 2026
% ============================================================================
% Ajouté à la demande de l'auteur (07/11/2026) : ultime chapitre qui ouvre le
% cours sur la recherche, avec des axes et des exemples de sujets.
% Sources : A0 « Vers l'auditologie » (hal-05701007) ; A1 « Modelling
% Epistemic Uncertainty in IS Audit » (hal-05701270) ; INF3226 C30 et C31.
% Règle de l'auteur (06/11/2026) : seuls A0 et A1 (HAL) et des travaux
% publiés sont cités ; aucun travail non publié n'est annoncé.
% Axiomes cités avec la numérotation d'A0 (1 à 6), pas avec l'environnement
% auto-numéroté.
% ============================================================================

\chapter[Ouvrir sur la recherche : l'auditologie]{Ouvrir sur la Recherche : l'Auditologie}
\label{chap:recherche}
\chaptermark{Ouvrir sur la recherche}

\epigraph{%
    « Ce que l'organisation ignore de ses propres systèmes conditionne la
    qualité des décisions qu'elle peut prendre à leur sujet. »%
}{\textit{Vers l'auditologie}~\cite{minka2026auditologie}}

\niveauChapitre{Recherche}{L'ensemble du cours, en particulier les
Chapitres~\ref{chap:preparer} à~\ref{chap:constater}.}

\begin{boxobjectifs}
\begin{itemize}
    \item Définir l'auditologie et la mission auditologique.
    \item Énoncer les six axiomes de l'article A0 et relier chacun à une
          pratique du cours.
    \item Présenter les cinq piliers théoriques et les prédictions
          falsifiables qu'ils engendrent.
    \item Distinguer ce qui est établi par des travaux publiés, ce qui est
          déposé en prépublication, et ce qui reste à valider.
    \item Identifier des axes de recherche à partir de la pratique de
          l'audit public, et formuler un sujet de mémoire falsifiable.
\end{itemize}
\end{boxobjectifs}

\bigskip

% ============================================================================
\section*{Pourquoi ce chapitre}
% ============================================================================

Tout au long du cours, des choix de méthode ont été justifiés par une
recherche~: la réponse \NEK{}, la décomposition en huit niveaux, les mesures
$\mu$, $U$ et $H$, le quadrant, la comparaison des cycles. Ces outils
viennent d'un programme de recherche que l'auteur a nommé
\emph{auditologie}~\cite{minka2026auditologie}, et dont l'instrument de
réponse est étudié dans un second article~\cite{minka2026nek}. Ce chapitre
le présente pour lui-même, puis montre comment un auditeur des finances
publiques peut y contribuer, par un mémoire ou par sa pratique.

\begin{boiteRemarque}{Ce que ce chapitre est, et n'est pas}
Il expose un programme de recherche déposé en prépublication, non un savoir
établi. Les axiomes et les piliers forment un cadre~; les prédictions qui en
découlent restent à tester. C'est précisément là que les auditeurs publics,
qui disposent de missions réelles et répétées, peuvent apporter ce qui
manque~: des données longitudinales.
\end{boiteRemarque}

\separateur

% ============================================================================
\section{L'objet : la mission auditologique}
\label{sec:c15-objet}
% ============================================================================

L'article A0 part d'un constat~: l'audit des systèmes d'information dispose
de méthodes, de référentiels et d'une abondante littérature instrumentale,
mais d'aucun socle axiomatique ni d'aucun espace de mesure unifié pour son
objet central, la réduction de l'incertitude. Il propose de le constituer en
une science appliquée.

\begin{boxdef}
\termecle{Mission auditologique}\index{mission auditologique} (A0)~:
quintuplet
$\mathcal{M} = \langle \mathcal{S}, \mathcal{A}, \mathcal{I},
\mathcal{G}, \mathcal{T} \rangle$, où $\mathcal{S}$ est le système audité
(ensemble sociotechnique d'agents, d'artefacts, de flux et de dispositifs),
$\mathcal{A}$ l'architecture cognitive de l'auditeur ou de l'équipe,
$\mathcal{I}$ le canal d'information par lequel les preuves parviennent à
$\mathcal{A}$, $\mathcal{G}$ la structure stratégique de l'interaction
entre l'auditeur, l'audité et les tiers, et $\mathcal{T}$ la trajectoire
temporelle du système au fil des missions.

\medskip
\termecle{Auditologie}\index{auditologie}~: science appliquée de la
réduction de l'incertitude sur les missions auditologiques.
\end{boxdef}

\begin{tableaucours}{cL{3.6cm}L{8.4cm}}{%
    \tableauHeader{} &
    \tableauHeader{Composante} &
    \tableauHeader{Dans la mission MINÉQUIR}
}
$\mathcal{S}$ & Système audité & La DAF, SIDÉPENSE, ses 47 comptes, ses
procédures, le poste comptable \\
$\mathcal{A}$ & Architecture cognitive & L'équipe du CONSUPE, sa méthode,
ses papiers de travail \\
$\mathcal{I}$ & Canal d'information & Entretiens, pièces des trois strates,
extractions faites devant la vérificatrice \\
$\mathcal{G}$ & Structure stratégique & DAF, cellule informatique, DRH,
contrôle financier, vérificateurs, CDBF \\
$\mathcal{T}$ & Trajectoire & Mars~2026, novembre~2026, et les cycles
suivants \\
\end{tableaucours}
\vspace{-0.8em}
\begin{center}{\small\itshape Tableau~15.1. La mission MINÉQUIR comme
mission auditologique}\end{center}

\separateur

% ============================================================================
\section{Six axiomes}
\label{sec:c15-axiomes}
% ============================================================================

L'article A0 pose six axiomes, qui contraignent toute mission
auditologique~\cite{minka2026auditologie}. Le cours les a rencontrés
chapitre après chapitre~; leur numérotation est celle de l'article.

\begin{tableaucours}{cL{3.4cm}L{5.0cm}L{3.6cm}}{%
    \tableauHeader{N°} &
    \tableauHeader{Axiome} &
    \tableauHeader{Énoncé résumé} &
    \tableauHeader{Dans le cours}
}
1 & Observabilité partielle & L'auditeur n'accède qu'à une partie des états
du système, elle-même partiellement opaque & Partir du système~; la réponse
\NEK{} (Ch.~\ref{chap:preuve}) \\
2 & Cognition limitée & L'auditeur et l'audité sont des agents
cognitivement limités & Décomposition bornée à sept éléments
(Ch.~\ref{chap:preparer}) \\
3 & Mesurabilité & L'incertitude épistémique se mesure par l'entropie et
ses généralisations & $\mu$, $U$, $H$ (Ch.~\ref{chap:constater}) \\
4 & Interaction stratégique & Auditeur et audité ont des intérêts
distincts, dans un jeu à information imparfaite & Biais du choix forcé,
indice $\rho$~; indépendance \\
5 & Individuation temporelle & Le système évolue, et son état dépend des
missions antérieures & Suivi, divergence entre cycles \\
6 & Opposabilité & Toute conclusion doit pouvoir être défendue devant un
tiers légitime & Papiers de travail, contradictoire, observation en sept
éléments \\
\end{tableaucours}
\vspace{-0.8em}
\begin{center}{\small\itshape Tableau~15.2. Les six axiomes de
l'auditologie (d'après A0)}\end{center}

A0 écarte explicitement trois candidats~: l'\emph{indépendance} n'est pas un
axiome, mais une conséquence des axiomes~4 et~6~; la \emph{conformité à un
référentiel} est un choix de fonction objectif~; la \emph{matérialité} est
un paramètre dérivé de l'opposabilité.

\subsection{Quatre propositions}

A0 en tire quatre propositions, que ce cours a mises en pratique.
\begin{enumerate}
    \item \textbf{Modalité d'ignorance explicite} (axiomes~1 à~4)~: tout
          instrument de réponse valide inclut \NEK{}.
    \item \textbf{Décomposition hiérarchique} (axiomes~2 et~6)~: tout
          instrument couvrant un espace d'audit non trivial est organisé en
          niveaux, bornés par la mémoire de travail.
    \item \textbf{Métriques longitudinales} (axiome~5)~: les conclusions ne
          se testent que par l'observation de plusieurs cycles.
    \item \textbf{Réflexivité} (axiomes~4 et~6)~: l'audit fait partie de la
          trajectoire du système qu'il observe.
\end{enumerate}

\separateur

% ============================================================================
\section{Cinq piliers}
\label{sec:c15-piliers}
% ============================================================================

Les axiomes disent ce qu'est une mission~; les piliers apportent les théories
qui permettent de la modéliser. A0 en retient cinq, chacun avec une
prédiction falsifiable.

\begin{tableaucours}{L{2.6cm}L{3.6cm}L{6.0cm}}{%
    \tableauHeader{Pilier} &
    \tableauHeader{Ancrage} &
    \tableauHeader{Prédiction falsifiable}
}
1. Cognition limitée & Miller, Simon,
Kahneman~\cite{miller1956,simon1957,kahneman2011} & Les grilles de plus de
neuf dimensions simultanées produisent plus d'incohérences et d'oublis \\
2. Théorie de l'information & Shannon, Kullback et
Leibler~\cite{shannon1948,kullback1951} & Le diagnostic conformité et
entropie prédit mieux que le seul taux de conformité \\
3. Théorie des jeux & von~Neumann, Nash,
Harsanyi~\cite{vonneumann1944,nash1951,harsanyi1967} & Les instruments
compatibles avec les incitations réduisent la dissimulation \\
4. Systèmes dynamiques & Chaînes de Markov~\cite{norris1998} & Les
prévisions par matrice de transition entre niveaux de maturité battent la
persistance naïve \\
5. Ontologie critique-réaliste & Bhaskar, Simondon,
Latour~\cite{bhaskar1975,simondon1958,latour2005} & Les cadres qui tiennent
compte des personnes, des procédures et des artefacts font mieux appliquer
les recommandations \\
\end{tableaucours}
\vspace{-0.8em}
\begin{center}{\small\itshape Tableau~15.3. Les cinq piliers et leurs
prédictions (d'après A0)}\end{center}

\separateur

% ============================================================================
\section{La méthode : science du design et falsifiabilité}
\label{sec:c15-methode}
% ============================================================================

\subsection{Une science du design}

L'auditologie produit des artefacts~: instruments, métriques, grilles,
outils. A0 l'inscrit dans la recherche en science du design~: on construit
un artefact, on l'évalue dans son environnement, on l'appuie sur la base de
connaissances, et l'on recommence~\cite{hevner2004design}. Les outils de ce
cours (l'instrument \OUI{}~/ \NON{}~/ \NEK{}, le gabarit canonique, les
labs de NDAP SUKULU) sont de tels artefacts, et chaque promotion
d'auditeurs participe à leur évaluation.

\subsection{Un noyau dur et une ceinture protectrice}

A0 adopte le falsificationnisme sophistiqué de Lakatos~\cite{lakatos1970}.
Les axiomes et les piliers forment un \emph{noyau dur}, qui n'est pas
directement testé~; les hypothèses auxiliaires, les procédures de mesure et
les prédictions forment la \emph{ceinture protectrice}, qui l'est. La
réfutation d'une prédiction n'emporte pas le programme~: elle oblige à
réviser la ceinture.

\subsection{Ce que l'auditologie n'est pas}

A0 précise quatre refus~: l'auditologie n'est pas \emph{normative} (elle ne
prescrit pas le bon audit), ni \emph{critique} (elle n'invalide pas les
cadres existants), ni \emph{un référentiel de plus}, ni \emph{une idéologie}
(ses propositions sont falsifiables). C'est pourquoi le cours a enseigné
ensemble les normes du métier (Partie~II) et leur lecture auditologique
(Partie~III).

\separateur

% ============================================================================
\section{Établi, déposé, à valider}
\label{sec:c15-statut}
% ============================================================================

\begin{tableaucours}{L{3.0cm}L{5.0cm}L{4.2cm}}{%
    \tableauHeader{Statut} &
    \tableauHeader{Contenu} &
    \tableauHeader{Références}
}
Établi par des travaux publiés & Borne de la mémoire de travail, entropie
de Shannon, divergence de Kullback-Leibler, théorie des jeux, contrôle
interne & Miller, Shannon, Kullback, Nash, COSO \\
Déposé en prépublication & Axiomes, piliers, mission auditologique,
instrument \OUI{}~/ \NON{}~/ \NEK{}, mesures, quadrant & A0 et A1, sur
HAL~\cite{minka2026auditologie,minka2026nek} \\
À valider & Les prédictions des cinq piliers, sur des missions réelles et
répétées & Les recherches à venir, dont les vôtres \\
\end{tableaucours}
\vspace{-0.8em}
\begin{center}{\small\itshape Tableau~15.4. Le statut des travaux sur
lesquels repose le cours}\end{center}

\separateur

% ============================================================================
\section{Axes de recherche pour l'auditeur public}
\label{sec:c15-axes}
% ============================================================================

Les auditeurs des finances publiques disposent de ce qui manque le plus à la
recherche en audit~: des missions réelles, sur des systèmes réels, répétées
dans le temps. Huit axes s'ouvrent à partir de ce cours. Pour chacun, un
exemple de sujet de mémoire est donné, avec sa question falsifiable.

\begin{tableaucours}{cL{3.2cm}L{5.0cm}L{3.0cm}}{%
    \tableauHeader{Axe} &
    \tableauHeader{Thème} &
    \tableauHeader{Question de recherche} &
    \tableauHeader{Pilier, axiome}
}
1 & Le biais du choix forcé & De combien les grilles binaires des
inspections surestiment-elles la conformité~? & Pilier~3~; axiome~4 \\
2 & L'incertitude dans les rapports publics & Quelle part des conclusions
des rapports de contrôle repose sur des déclarations sans preuve~? &
Pilier~2~; axiome~1 \\
3 & La valeur prédictive du quadrant & Les domaines en conformité opaque
révèlent-ils plus d'irrégularités aux missions suivantes~? & Pilier~2~;
axiome~3 \\
4 & Trajectoires des recommandations & Le suivi des recommandations
suit-il une chaîne de Markov prévisible~? & Pilier~4~; axiome~5 \\
5 & L'usage stratégique de \NEK{} & Les audités répondent-ils davantage
\NEK{} en entretien formel qu'en auto-évaluation~? & Pilier~3~; axiome~4 \\
6 & La décomposition et l'erreur & Les grilles qui dépassent la borne de
sept produisent-elles plus d'oublis~? & Pilier~1~; axiome~2 \\
7 & Contrôle interne des applications de l'État & Quelles déficiences
reviennent dans les applications financières publiques, et pourquoi~? &
Pilier~5~; axiome~1 \\
8 & L'opposabilité de la preuve électronique & Quelles conditions rendent
une extraction recevable devant le CDBF et la Chambre des Comptes~? &
Axiome~6 \\
\end{tableaucours}
\vspace{-0.8em}
\begin{center}{\small\itshape Tableau~15.5. Huit axes de recherche ouverts
par le cours}\end{center}

\subsection{Exemples de sujets}

\begin{boxscenario}{Axe 1. Mesurer le biais des grilles binaires}
\textbf{Sujet.} \emph{Le biais du choix forcé dans les questionnaires de
contrôle interne des administrations financières camerounaises~: mesure par
l'indice $\rho$.}

\textbf{Hypothèse falsifiable.} Sur les mêmes questions et les mêmes
répondants, la part de \OUI{} d'une grille binaire dépasse d'au moins dix
points celle de l'instrument trichotomique.

\textbf{Données.} Deux passages d'un même questionnaire (binaire puis
trichotomique) dans trois à cinq services, avec l'accord de la hiérarchie.

\textbf{Méthode.} Calcul de $\rho$ par question et par domaine~; test
statistique de l'écart~; analyse des questions où le biais est le plus fort.

\textbf{Réfutation.} Un $\rho$ proche de zéro sur l'ensemble des services.
\end{boxscenario}

\begin{boxscenario}{Axe 3. La valeur prédictive du quadrant}
\textbf{Sujet.} \emph{Conformité opaque et irrégularités ultérieures~: étude
longitudinale de missions de vérification.}

\textbf{Hypothèse falsifiable.} Les domaines classés en conformité opaque à
une mission présentent, à la mission suivante, plus d'irrégularités établies
que les domaines en conformité claire.

\textbf{Données.} Les mesures $\mu$, $U$, $H$ de domaines audités sur deux
cycles au moins, et les constats de la mission suivante.

\textbf{Méthode.} Comparaison des taux d'irrégularités par région~; contrôle
de la taille des domaines.

\textbf{Réfutation.} Aucune différence entre régions, ou une différence
inverse.
\end{boxscenario}

\begin{boxscenario}{Axe 4. Trajectoires des recommandations}
\textbf{Sujet.} \emph{Le suivi des recommandations de contrôle comme
processus de Markov~: estimation d'une matrice de transition des niveaux de
maturité.}

\textbf{Hypothèse falsifiable.} Une matrice de transition estimée sur trois
cycles prévoit le niveau de maturité au cycle suivant mieux que la simple
reconduction du niveau actuel.

\textbf{Données.} Les niveaux de maturité (échelle du
Chapitre~\ref{chap:eval-sci}) de plusieurs domaines, sur au moins quatre
cycles.

\textbf{Méthode.} Estimation de la matrice, prévision, comparaison à la
persistance naïve.

\textbf{Réfutation.} Aucun avantage prédictif de la matrice.
\end{boxscenario}

\begin{boxscenario}{Axe 8. L'opposabilité de la preuve électronique}
\textbf{Sujet.} \emph{Les conditions de recevabilité des extractions
informatiques dans les procédures du CDBF et de la Chambre des Comptes.}

\textbf{Question.} Quelles garanties (origine, exhaustivité, exactitude,
intégrité) sont effectivement exigées, et lesquelles manquent le plus
souvent~?

\textbf{Données.} Un échantillon de dossiers, anonymisés, avec l'autorisation
des institutions~; entretiens avec des rapporteurs.

\textbf{Méthode.} Grille d'analyse fondée sur les quatre garanties du
Chapitre~\ref{chap:preuve}~; étude des cas où une pièce a été écartée.
\end{boxscenario}

\separateur

% ============================================================================
\section{Conduire un mémoire de recherche appliquée}
\label{sec:c15-memoire}
% ============================================================================

\begin{boxprocedure}
\textbf{Les étapes d'un mémoire en auditologie appliquée}~:
\begin{enumerate}
    \item partir d'une difficulté réelle de mission, et la rattacher à un
          axiome ou à un pilier~;
    \item formuler une hypothèse \emph{falsifiable}~: dire à l'avance quel
          résultat la réfuterait~;
    \item obtenir l'autorisation écrite de la hiérarchie et des entités
          concernées~;
    \item construire le protocole~: instrument, échantillon, mesures
          ($\mu$, $U$, $H$, $\rho$, $D_{\mathrm{KL}}$), seuils déclarés~;
    \item anonymiser les données personnelles et protéger les extractions~;
    \item analyser, puis séparer ce qui est établi, ce qui est probable et ce
          qui reste inconnu~;
    \item discuter les limites, et proposer la révision de la ceinture
          protectrice si la prédiction est réfutée.
\end{enumerate}
\end{boxprocedure}

\begin{boxalert}
\textbf{Violation typique.} Formuler un sujet du type «~l'importance de
l'audit des SI dans l'administration camerounaise~».

\smallskip
\textbf{Règle~:} un sujet de recherche pose une question à laquelle des
données peuvent répondre non. «~Importance~» ne se réfute pas~; «~les
domaines en conformité opaque révèlent plus d'irrégularités au cycle
suivant~» se réfute.
\end{boxalert}

\begin{boxmetier}{Vérificateur}
Vos missions répétées sur les mêmes entités sont des données longitudinales
rares. Conservez, d'un cycle à l'autre, les mêmes questions et les mêmes
fonctions répondantes~: vous rendez possibles les axes~3 et~4.
\end{boxmetier}

\begin{boxmetier}{Auditeur interne}
Mesurer $\rho$ sur vos propres questionnaires est un premier travail de
recherche, utile dès la première année.
\end{boxmetier}

\separateur

\begin{boxresume}
\begin{itemize}
    \item \textbf{L'auditologie} est la science appliquée de la réduction de
          l'incertitude~; son objet est la mission
          $\langle \mathcal{S}, \mathcal{A}, \mathcal{I}, \mathcal{G},
          \mathcal{T} \rangle$.
    \item \textbf{Six axiomes}~: observabilité partielle, cognition limitée,
          mesurabilité, interaction stratégique, individuation temporelle,
          opposabilité.
    \item \textbf{Quatre propositions} fondent les pratiques du cours~:
          \NEK{}, décomposition, mesures longitudinales, réflexivité.
    \item \textbf{Cinq piliers}, chacun avec une prédiction falsifiable.
    \item \textbf{Une science du design, à la Lakatos}~: un noyau dur, une
          ceinture protectrice qui se teste.
    \item \textbf{Un statut clair}~: A0 et A1 sont des prépublications~; les
          prédictions restent à valider.
    \item \textbf{Huit axes} s'ouvrent aux auditeurs publics~; un bon sujet
          dit à l'avance ce qui le réfuterait.
\end{itemize}
\end{boxresume}

\begin{boxexo}[title={\ding{45}\quad Exercice~15.1. Axiomes et pratiques \niveaubase}]
Pour chacune des pratiques suivantes, citez l'axiome d'A0 qui la fonde~:
(a) coter \NEK{} une pièce non encore remise~; (b) borner chaque niveau de
la décomposition à sept éléments~; (c) comparer deux cycles par la
divergence~; (d) soumettre le rapport provisoire à l'entité.
\end{boxexo}

\begin{boxexo}[title={\ding{45}\quad Exercice~15.2. Rendre un sujet falsifiable \niveaumoyen}]
Réécrivez les sujets suivants pour qu'ils deviennent falsifiables, en
précisant l'hypothèse, les données et la condition de réfutation~: (a)
«~Le rôle du CONSUPE dans la bonne gouvernance numérique~»~; (b) «~Les
enjeux de la dématérialisation des recettes~».
\end{boxexo}

\begin{boxexo}[title={\ding{45}\quad Exercice~15.3. Un protocole \niveauexpert}]
Choisissez l'un des huit axes du Tableau~15.5 et rédigez le protocole
complet d'un mémoire~: question, hypothèse, données et autorisations,
instrument, mesures, seuils, analyse, condition de réfutation, risques
éthiques et parade.
\end{boxexo}

\begin{boxinfo}
\textbf{Pour finir.} Le cours s'arrête ici~; le métier, et peut-être la
recherche, commencent. Chaque mission refera les mêmes gestes~: partir du
système, écrire ses critères avant de tester, coter honnêtement, y compris
\NEK{}, conclure sur ce que l'on a prouvé, et mesurer, d'un cycle à
l'autre, ce que l'on a appris.
\end{boxinfo}
